\documentclass[12pt, a4paper]{article}
\usepackage[T1]{fontenc}
\usepackage[utf8]{inputenc}
\usepackage{lmodern}
\usepackage{microtype}
\usepackage{amsmath, amssymb}
\usepackage{booktabs}
\usepackage{graphicx}
\usepackage{caption}
\usepackage[style=authoryear, backend=biber]{biblatex}
\addbibresource{references.bib}
\usepackage[margin=2.5cm]{geometry}
\usepackage{setspace}
\usepackage{xcolor}
\usepackage{hyperref}
\usepackage{threeparttable}

\onehalfspacing

\newcommand{\todo}[1]{\textcolor{red}{\textbf{[TODO: #1]}}}

\title{Technology fields, automation and task reinstatement\\
\large Working paper --- not for circulation}
\author{Joakim Storck and Jonatan Andersson}
\date{Draft: \today}

\begin{document}
\maketitle

\begin{abstract}
\noindent
We use the geometric task-space framework of \textcite{storck2026geometry}
to study how the reach of a technology determines whether task creation
benefits or harms workers in different parts of the labour market.
The framework places occupations as task bundles in a two-dimensional
task space calibrated from occupational embeddings. The market price of
skill is a field over this space, directional rather than uniform, so
that depth is rewarded in the socio-cognitive directions and not in the
service and manual directions. Automation removes the high-priced tasks
from a bundle; reinstatement returns a fraction $\gamma$ of the displaced
mass as new tasks, which attach to an occupation only where the priced
capability they require is already held. Occupations are therefore
redefined rather than abolished, and their wages move with the priced
content of their bundles. New tasks that no occupation can attach remain
unbound, and their concentration marks where new occupations may emerge.

We calibrate the technology parameters against task-level AI exposure
scores from \textcite{gmyrek2025} and test the model's central spatial
prediction against Bureau of Labor Statistics occupational employment
data for 741 occupations over 2019--2024. A difference-in-differences
specification exploiting the ChatGPT launch of November 2022 shows
that occupations at the gradient boundary of the augmentation field
experienced significantly higher employment growth in the post-2022
period, while occupations in the field's core experienced lower growth.
The boundary zone effect is robust across six specifications and holds
in an event study with clean pre-trends.

The model also reveals a structural threshold at $z_K \approx
0.38$--$0.40$. Below this threshold, reinstatement concentrates new
tasks in the cognitive and investigative sectors and harms manual
workers through general equilibrium redistribution. Above it, the
technology's gradient boundary reaches the manual arc and reinstatement
benefits manual workers in employment, though the directional price of
skill mutes the wage benefit in those directions. This spatial
condition, the reach of the technology relative to the distance between
the technology centre and the manual arc, cannot be derived from a
one-dimensional task-space model.
\end{abstract}

\tableofcontents
\newpage

%% ============================================================
\section{Introduction}
\label{sec:intro}
%% ============================================================

The division of labor segments the labor market into specialized
domains. Workers accumulate human capital specific to their occupation
or sector, and technologies are adopted unevenly across the task
distribution. When a technology displaces workers from one set of
tasks, the consequences for wages and employment depend on whether
new tasks are created to absorb the displaced labor: the reinstatement
mechanism of \textcite{acemoglu2019automation}. In the canonical
one-dimensional framework, reinstatement unambiguously raises the
labor share because new tasks are by construction assigned to labor
at the boundary of the technology's operating domain
\parencite{acemoglu2018race}. The question
of which workers benefit is not well posed in a model where labor
enters as two or three aggregate skill groups distinguished only by
their position on a single task index.

This paper develops a spatial generalization of the task-based
framework and tests its central prediction against occupational
employment data. The framework places occupations as task bundles in
a two-dimensional task space calibrated from semantic embeddings of
O*NET task statements \parencite{storck2026geometry}. The two
dimensions capture domain orientation and specialization depth,
placing occupations in a continuous space that reflects the actual
content of work rather than predefined skill categories. The market
price of skill is a field over this space, directional rather than
uniform. Technologies are represented as augmentation fields: Gaussian
productivity enhancements centered at a point in task space with a
radius that determines their reach across occupational domains.

The spatial structure generates a precise prediction about where new
tasks appear. The gradient of the augmentation field is maximized at
distance $z_K$ from the technology center: at the boundary between
the core, where the technology dominates and the marginal value of
complementary human activity is low, and the periphery, where the
technology is irrelevant. New tasks that validate, coordinate, and
integrate automated outputs emerge precisely in this boundary zone.
Whether reinstatement benefits workers therefore depends not on the
reinstatement elasticity alone, but on where the gradient boundary
falls relative to the spatial ownership structure of task space, and
on the price the market attaches to work in that direction. A
broad, inexpensive technology centered in the cognitive arc can
generate task creation that harms manual workers through general
equilibrium redistribution rather than offsetting their displacement.

We test this prediction using Bureau of Labor Statistics Occupational
Employment and Wage Statistics (OEWS) data for 741 occupations over
2019--2024, combined with task-level AI exposure scores from
\textcite{gmyrek2025} mapped onto the occupational geometry. The
treatment variable is not binary but continuous: for each occupation
we compute both its exposure to the augmentation field and its position
relative to the field's gradient. A difference-in-differences
specification exploits the ChatGPT launch of November 2022 as an
exogenous shift in the technology's reach. Occupations at the gradient
boundary show significantly higher employment growth in the post-2022
period relative to pre-trends, while occupations in the field's core
show lower growth. The boundary zone effect is robust across six
specifications and holds in an event study with clean pre-trends.

The paper proceeds as follows. Section~\ref{sec:geometry} introduces
the spatial occupational geometry. Section~\ref{sec:model} presents
the task-based production framework and its spatial extension,
including the price of skill, the augmentation field, displacement,
and displacement-anchored reinstatement through a bundle operator.
Section~\ref{sec:calibration} describes the calibration of technology
parameters against empirical AI exposure data. Section~\ref{sec:cases}
illustrates the model with two scenarios drawn from the
Acemoglu--Restrepo taxonomy. Section~\ref{sec:empirics} presents the
empirical analysis. Section~\ref{sec:discuss} discusses implications
and limitations.

%% ============================================================
\section{The Spatial Occupational Geometry}
\label{sec:geometry}
%% ============================================================

\todo{Kort presentation av geometrin från Paper 1: taskspace,
vinkeldimension (domän), radiedimension (specialisering),
occupational centroids och kompetensradier. En figur (taskspace-kartan).
Ca 1--1.5 sidor.}

The value of this representation can be shown through its ability to illustrate successive waves of labor-replacing technology. 
The routine-task automation studied by \textcite{autor2003} targeted occupations near the Technical-Physical pole, contributing to the secular decline of middle-wage manufacturing employment \parencite{acemoglu2020, graetz2018, devries2020}. 
Generative AI capabilities \parencite{hampole2025, eloundou2024} extend into parts of the Social-Cognitive region previously considered safe, raising concerns about the exposure of college-educated work. 
Each technology can be characterized by the region of the task space it affects, and each occupation by its position relative to that region. 
This is a question the geometry is built to address but that a partitioned or one-dimensional model cannot express.

%% ============================================================
\section{Model}
\label{sec:model}
%% ============================================================

We embed the geometry into a production framework that builds on the
canonical task-based model of \textcite{acemoglu2019automation},
extending it by replacing abstract factor categories with a
continuous task space and by letting technology redraw the task
bundles that define occupations.

\subsection{Overview}
\label{sec:overview}

The model has two layers. A \emph{task layer} is a set of fields over
the task disk. It carries production, the technology, displacement,
reinstatement, and the creation of unbound tasks. A \emph{worker
layer} carries employment, mobility, and labor supply. The two layers
are linked by the bundles that define occupations, and it is these
bundles that automation and reinstatement rewrite.

The task layer is family-free. Job families are an external
classification by industrial context, and the geometry of
\textcite{storck2026geometry} is built without such categories. We work
at the level of tasks, the level at which the geometry is constructed
and at which technology exposure is measured. Occupations enter as
bundles of tasks, not as primitives.

Wages are pinned by the price of skill and the task bundles alone, and
the model carries no occupation-specific wage parameter. The price of
skill is a field over the disk, identified by
\textcite{storck2026geometry} and carried here as the price of labor at
each task location.

\subsection{The Price of Skill}
\label{sec:price}

\textcite{storck2026geometry} shows that the labor market does not
value position in task space isotropically. The return to radial depth
is positive toward the analytical directions and absent toward the
service and manual directions, and what the market prices is the
socio-cognitive cluster of capabilities concentrated in the northern
half. We summarize this as a price field $\Pi(\mathbf{r})$ over the
disk, the price the market pays for a unit of work performed at task
location $\mathbf{r} = (\xi, \chi)$. It is the wage field of
\textcite{storck2026geometry}, a first-harmonic field on the circle
whose radial term is modulated by direction,
\begin{equation}
  \ln \Pi(\mathbf{r})
  \;=\;
  m_0 + m_1 \cos\xi + m_2 \sin\xi
  \;+\; \chi\,\bigl(m_3 + m_4 \cos\xi + m_5 \sin\xi\bigr).
  \label{eq:price-field}
\end{equation}
The direction-modulated radial term is what makes depth rewarded in
the north and unrewarded in the south. $\Pi$ is estimated once, on the
cross-section of occupations, and treated as fixed structure here. It
is the price of labor at a task location.

Productivity is isotropic. \textcite{storck2026geometry} does not
find that workers are more productive in the north. It finds that the
market pays more for the capabilities concentrated there. Productivity
belongs to the worker and is symmetric; the price of that work is
directional. The directional structure therefore lives in $\Pi$, not
in a productivity kernel.

\subsection{Tasks, Occupations, and Wages}
\label{sec:tasks-occ}

A task is a unit of work activity that yields a unit of output. Because
one task is one unit of output, no productivity magnitude varies across
tasks. What varies over the disk is how much labor is present and what
that labor is paid.

An occupation $o$ is a bundle: a weight measure $b_o(\mathbf{r})$ over
task space, normalized to one, with the importance weights of
\textcite{storck2026geometry}. Its centroid
$\mu_o = \int \mathbf{r}\, b_o(\mathbf{r})\, d\mathbf{r}$ is a summary
of the bundle, not the object itself. The occupation's wage is the
price of its bundle,
\begin{equation}
  w_o \;=\; \int \Pi(\mathbf{r})\, b_o(\mathbf{r})\, d\mathbf{r},
  \label{eq:wage-bundle}
\end{equation}
the average price of skill over the tasks the occupation performs. The
wage is therefore not a property of the occupation but an aggregate of
the price field over its current bundle, and it moves when the bundle
is rewritten. Labor density on the disk is
\begin{equation}
  n(\mathbf{r}) \;=\; \sum_o L_o\, b_o(\mathbf{r}),
  \label{eq:labor-density}
\end{equation}
where $L_o$ is employment in occupation $o$, the employment-weighted
task coverage of the disk.

\subsection{Technology and the Operated Regime}
\label{sec:regime}

A technology $K$ is a field $\phi_K(\mathbf{r})$, a
Gaussian centered at $p_K$ with reach $z_K$,
\begin{equation}
  \phi_K(\mathbf{r})
  \;=\;
  A_K \exp\!\left(-\frac{1}{2}
  \left(\frac{\|\mathbf{r} - p_K\|}{z_K}\right)^2\right),
  \label{eq:phi-K}
\end{equation}
the effectiveness with which capital can do the work at $\mathbf{r}$.
Capital does not perform a task autonomously. At a location the worker
remains and the output is unchanged; what the technology changes is the
share of the location's micro-tasks that capital bears. The position
$p_K$ fixes where in task space the technology bites, and the reach
$z_K$ fixes how far it extends: a narrow $z_K$ is a specialized
technology, a broad $z_K$ a general-purpose one. Producing the crossed
micro-tasks more cheaply with capital is the productivity effect of
automation in the sense of \textcite{acemoglu2018race}, here made
spatial.

A technology also has a character. The same effectiveness can replace
the worker or assist her, and these have different consequences. We
carry the character as a scalar $s_K \in [0,1]$, a fourth attribute of
the technology beside $p_K$, $z_K$, and $A_K$. At $s_K = 1$ the
technology is purely replacing, at $s_K = 0$ purely augmenting. The
effectiveness splits into a replacing part $s_K\,\phi_K$ and an
augmenting part $(1 - s_K)\,\phi_K$, which act on the same location
through different channels. The augmenting part assists rather than
replaces; we treat it in Section~\ref{sec:character}.

Whether capital bears a micro-task is a question of price, and only the
replacing part competes with the worker on it. Producing the task with
labor costs $\Pi(\mathbf{r})$, the price of the work. Producing it with
the replacing part of capital costs the rental $R$ against that part's
effectiveness, an effective cost $R/(s_K\,\phi_K(\mathbf{r}))$. Capital
is the cheaper producer when
\begin{equation}
  s_K\,\phi_K(\mathbf{r}) \;>\; \frac{R}{\Pi(\mathbf{r})}.
  \label{eq:regime-shift}
\end{equation}
The comparison is between the rental of capital and the price of labor
at the location, weighted by the technology's effectiveness there. The
threshold $R/\Pi(\mathbf{r})$ is lowest where $\Pi(\mathbf{r})$ is
highest, so capital is profitable first at the most highly priced
locations, and a broad, inexpensive technology can sustain it well
beyond its core. Capital moves first against the dear work. This is the
spatial counterpart of the threshold in \textcite{acemoglu2018race}:
where they automate the lowest task on a single scale, the field takes
the highest-priced positions wherever they lie on the disk.

A location $\mathbf{r}$ is not a single task. It is a neighborhood
holding a continuum of micro-tasks, and the effectiveness $\phi_K$
varies across them, so \eqref{eq:regime-shift} is not true or false for
the whole cell at once. It holds for a fraction of the micro-tasks and
fails for the rest. We write that fraction as $a(\mathbf{r})$, the
degree to which capital bears the work at $\mathbf{r}$,
\begin{equation}
  a(\mathbf{r})
  \;=\;
  \frac{1}{1 + \exp\!\bigl(-(s_K\,\phi_K(\mathbf{r}) - R/\Pi(\mathbf{r}))/\tau\bigr)},
  \label{eq:soft-switch}
\end{equation}
where $\tau$ sets the width of the margin, the dispersion of
\eqref{eq:regime-shift} across the micro-tasks of the cell. As
$\tau \to 0$ the cell collapses to a single task and $a(\mathbf{r})$
becomes a step, zero below the threshold and one above it; in that
limit the operated region is a sharp partition of task space and the
margin is a boundary in space. For $\tau > 0$ the margin lies inside
each cell, and $a(\mathbf{r})$ is the share of its micro-tasks on the
operated side of \eqref{eq:regime-shift}.

The replacing part sets the takeover; the augmenting part does not enter
it. At $s_K = 1$ the threshold \eqref{eq:regime-shift} and the operated
share \eqref{eq:soft-switch} are those of pure replacement, and the
displacement, wage, and density results that follow hold unchanged. The
replacement model is the case $s_K = 1$.

Displacement is then a loss of priced contribution, not a loss of the
task. The worker remains at $\mathbf{r}$ and the output is unchanged;
what shrinks is the share of that output the worker is paid for. The
priced contribution of the cell is $\Pi(\mathbf{r})$ scaled by the
retained share $1 - a(\mathbf{r})$. The mass displaced from occupation
$o$ is the operated share integrated over its own bundle,
\begin{equation}
  D_o \;=\; \int b_o(\mathbf{r})\, a(\mathbf{r})\, d\mathbf{r},
  \label{eq:displaced}
\end{equation}
and its retained wage is
$\int \Pi(\mathbf{r})\, b_o(\mathbf{r})\, (1 - a(\mathbf{r}))\,
d\mathbf{r}$. The aggregate displaced labor mass is
$\Delta\Gamma^D = \sum_o L_o D_o$. An occupation can carry substantial
$D_o$ without any region having become operated as a whole, because the
displacement is spread across its bundle rather than concentrated in a
zone it owns. This is the difference between an occupation disappearing
and an occupation that remains but is paid for less of what it does.

Because $a(\mathbf{r})$ varies within the bundle, the margin cuts
inside the occupation rather than between regions. Some of an
occupation's own tasks are borne by capital while others remain with
the worker. This is the restructuring of work within an occupation, the
unbundling and rebundling of \textcite{brynjolfsson2018mlworkforce},
expressed as a field. The reach of the technology, $p_K$ and $z_K$,
together with the price field $\Pi$, determines which tasks within each
bundle cross the margin.

\subsection{Reinstatement: The Bundle Operator}
\label{sec:reinstatement}

Reinstatement returns a fraction $\gamma$ of the displaced mass as new
tasks, and the new tasks must attach to occupations. Free tasks have
no role in the model. We therefore need a rule for where new tasks
appear and which occupations can take them.

\paragraph{Where they appear.}
The gradient of the augmentation field,
\begin{equation}
  \|\nabla \phi_K(\mathbf{r})\|
  \;=\;
  \frac{1}{z_K^2}\,\|\mathbf{r} - p_K\| \cdot \phi_K(\mathbf{r}),
  \label{eq:gradient-phi}
\end{equation}
peaks at distance $z_K$ from $p_K$. New tasks are seeded on this
boundary ring. With $\hat{g}(\mathbf{r}) = \|\nabla\phi_K(\mathbf{r})\|
/ \int \|\nabla\phi_K\|$ the normalized ring density and $M = \gamma\,
\Delta\Gamma^D$ the seeded mass, the seeded density is
$s(\mathbf{r}) = M\,\hat{g}(\mathbf{r})$.

\paragraph{Which occupations can take them.}
A new task at $\mathbf{r}$ requires a profile of capabilities, read from
the projected skill and ability profiles of
\textcite{storck2026geometry} evaluated at $\mathbf{r}$. An occupation
can attach the task readily when it already holds those capabilities,
and only with difficulty when the task requires a capability the
occupation lacks. The relevant quantity is the deficit, the positive
shortfall, not a symmetric distance:
\begin{equation}
  \delta_o(\mathbf{r})
  \;=\;
  \sum_k v_k \,\max\!\bigl(q_k(\mathbf{r}) - q_{o,k},\, 0\bigr),
  \label{eq:deficit}
\end{equation}
where $q_k(\mathbf{r})$ is the level of capability cluster $k$ required
at $\mathbf{r}$, $q_{o,k}$ is the occupation's level on cluster $k$, and
$v_k$ is the weight of cluster $k$. Attaching tasks that use less of the
occupation's capabilities is free. Attaching tasks that demand more, or
a different kind, is costly.

The weights $v_k$ are the cluster wage returns from the mediation
analysis of \textcite{storck2026geometry}: large for the
socio-cognitive skill cluster S1, small for the technical cluster S2,
near zero for the ability clusters A1 and A2. A single set of weights
does two jobs. It measures how expensive a capability is to acquire,
and how much the market pays for it. New tasks attach freely only where
the deficit is small, that is, where the priced capability content does
not exceed what the occupation already holds. Reinstatement can
therefore broaden a bundle but cannot lift its priced content. The wage
cannot jump. A large jump would require a capability the occupation does
not have, and that is not reinstatement within an occupation but the
seed of a new one.

In the one-dimensional framework, comparative advantage is a monotone
schedule: labor is strongest in the highest-indexed tasks, and new
tasks are added there by construction (Assumption~1 of
\textcite{acemoglu2018race}). The deficit gate is the geometric
counterpart. It does not posit a single direction in which labor is
strong. It lets a new task attach where the occupation already holds
the priced capability the task requires, wherever that lies on the
disk.

\paragraph{Attachment and allocation.}
The occupation's readiness to attach a task at $\mathbf{r}$ is
\begin{equation}
  e_o(\mathbf{r}) \;=\; \exp\!\bigl(-\delta_o(\mathbf{r})/\ell\bigr),
  \label{eq:readiness}
\end{equation}
between zero and one, where $\ell$ is a single global scale. The total
attachment capacity at $\mathbf{r}$ is $C(\mathbf{r}) = \sum_o L_o\,
e_o(\mathbf{r})$, the labor present weighted by its readiness. A
saturating function $\Phi(C) = C/(1 + C)$ gives the bound fraction of
the seeded mass. Occupation $o$ receives the reinstated density
\begin{equation}
  \iota_o(\mathbf{r})
  \;=\;
  s(\mathbf{r})\,\Phi(C(\mathbf{r}))\,
  \frac{L_o\, e_o(\mathbf{r})}{C(\mathbf{r})},
  \label{eq:reinstated}
\end{equation}
its share of the bound mass in proportion to its readiness.

\paragraph{The bundle operator.}
The post-technology bundle removes displaced weight and adds reinstated
weight,
\begin{equation}
  b_o^{\text{post}}(\mathbf{r})
  \;=\;
  b_o(\mathbf{r})\,\bigl(1 - a(\mathbf{r})\bigr)
  \;+\;
  \frac{\iota_o(\mathbf{r})}{L_o}.
  \label{eq:bundle-operator}
\end{equation}
The same field $\phi_K$ drives both terms. It deactivates through
$a(\mathbf{r})$, high near the core, and seeds through
$\hat{g}(\mathbf{r})$, high on the ring. The bundle is rewritten
continuously as the technology intensifies. This is the geometric form
of occupational redefinition \parencite{brynjolfsson2018mlworkforce}:
the occupation does not vanish, its bundle is redrawn.

\paragraph{Wages and the Turing trap.}
The wage change is the price-weighted change in the bundle,
\begin{equation}
  \Delta w_o
  \;=\;
  \int \Pi(\mathbf{r})\,
  \bigl[b_o^{\text{post}}(\mathbf{r}) - b_o(\mathbf{r})\bigr]\,
  d\mathbf{r}
  \;=\;
  \int \Pi(\mathbf{r})\,
  \Bigl[-b_o(\mathbf{r})\,a(\mathbf{r})
  + \tfrac{\iota_o(\mathbf{r})}{L_o}\Bigr]\, d\mathbf{r}.
  \label{eq:wage-change}
\end{equation}
The first term is negative and falls on high-$\Pi$ tasks; the second
adds tasks that, by the deficit gate, are mostly low-$\Pi$. Default
reinstatement therefore carries a downward wage bias: it strips priced
work and refills with unpriced work. The sign of $\Delta w_o$ is the
contest between where $a(\mathbf{r})$ removes weight and where
$\iota_o$ adds it, read against the price gradient $\nabla\Pi$. In the
directions where depth is priced, the strip falls on high-$\Pi$ tasks
and the refill cannot replace them.

\paragraph{Density.}
Summing the bundle change over occupations gives the evolved task
density,
\begin{equation}
  \rho_1(\mathbf{r}) - \rho_0(\mathbf{r})
  \;=\;
  -\sum_o L_o\, b_o(\mathbf{r})\, a(\mathbf{r})
  \;+\;
  \sum_o \iota_o(\mathbf{r})
  \;+\;
  u(\mathbf{r}),
  \label{eq:rho-evolution}
\end{equation}
with $u(\mathbf{r})$ defined in Section~\ref{sec:unbound}. The added
density integrates to $\gamma\,\Delta\Gamma^D$, so $\gamma$ keeps its
interpretation as the fraction of displaced mass returned. Setting
$\gamma = 0$ recovers pure automation; $\gamma = 1$ implies full
reinstatement.

\subsection{Unbound Tasks and Candidate Occupations}
\label{sec:unbound}

Not all seeded mass attaches. Where every occupation has a large
deficit, $C(\mathbf{r})$ is small, $\Phi(C(\mathbf{r}))$ is near zero,
and the seeded mass stays unbound:
\begin{equation}
  u(\mathbf{r})
  \;=\;
  s(\mathbf{r})\,\bigl[1 - \Phi(C(\mathbf{r}))\bigr].
  \label{eq:unbound}
\end{equation}
$u$ is the density of new tasks that no existing occupation can perform.
It is an output of the model, not an agent. The bound mass
$\sum_o \iota_o$ and the unbound mass $u$ together integrate to the
seeded $\gamma\,\Delta\Gamma^D$, so the closed-occupation accounting
holds.

The map of $u$ answers two distinct questions. Where unbound tasks are
created is $u$ itself, the ring regions that are active but that no
bundle can absorb. Where new occupations may emerge is where $u$
concentrates, a connected region carrying enough mass to support an
occupation, not merely a high value at a point. A diffuse $u$ along the
ring is friction. A clump is a candidate.

The price field tells the two apart. Unbound mass in a high-$\Pi$ region
is latent skilled work awaiting a bearer, the kind of region from which
qualified new occupations have historically grown. Unbound mass in a
low-$\Pi$ region is work no occupation wants and no worker is paid well
for, closer to precarious friction than to a new occupation. We color
candidate regions by $\Pi$ and read opportunity against spillage. We do
not instantiate new occupations: the model shows where one may appear,
not that one does.

The same deficit that gates attachment draws the boundary between two
literatures. New tasks are absorbed by existing occupations where the
priced deficit is small, which is redefinition in the sense of
\textcite{brynjolfsson2018mlworkforce}. They remain unbound where the
deficit is large for every occupation, which is the new work of
\textcite{acemoglu2019automation} that no existing occupation owns. One
operator decides which, through the deficit.

\subsection{Employment and Mobility}
\label{sec:equilibrium}

The worker layer sets employment $L_o$ and closes the model. We solve it
as a static spatial-sorting equilibrium and read the effect of a
technology as the comparative static between two such equilibria, one
with the original bundles and one with the bundles rewritten by
\eqref{eq:bundle-operator}. A law of motion in time can be added on top
of this rest point; we do not need it for a statement about the new
equilibrium.

\paragraph{Two factors share each place.}
A location is produced by human and capital input as perfect
substitutes, in the manner of \textcite{acemoglu2018race}. With density
$n(\mathbf{r}) = \sum_o L_o\, b_o(\mathbf{r})$ from
\eqref{eq:labor-density}, the human effective work at $\mathbf{r}$ is
$n(\mathbf{r})(1 - a(\mathbf{r}))$ and the capital effective work is
$n(\mathbf{r})\,a(\mathbf{r})$, so the total task-work is
$n(\mathbf{r})$ regardless of the split. Output value is
\begin{equation}
  V(\mathbf{r})
  \;=\;
  \Pi(\mathbf{r})\, n(\mathbf{r})^{\beta},
  \qquad \beta \in (0,1),
  \label{eq:place-output}
\end{equation}
concave in task-work through the congestion exponent $\beta$, and
unchanged by who performs the work. We hold the output value of a place
fixed under the human-capital split. This keeps the framework partial:
it carries no channel for goods prices or growth, which would require
modeling the wider economy, and we leave that to future work.
Congestion acts on $n$, the work present; the operated share $a$ enters
separately, as the part of the work the human no longer performs. Because
the bundle $b_o$ is normalized, a broad task radius spreads it thin and
contributes little to $n(\mathbf{r})$, so a concentrated occupation
presses harder on the places it occupies than a diffuse one does.

\paragraph{Returns to the two factors.}
Human and capital are paid the marginal product of task-work,
$\beta\,\Pi(\mathbf{r})\, n(\mathbf{r})^{\beta-1}$. The return to a
worker is that marginal product on the share she still performs,
\begin{equation}
  \rho(\mathbf{r})
  \;=\;
  \beta\,\Pi(\mathbf{r})\,\bigl(1 - a(\mathbf{r})\bigr)\,
  n(\mathbf{r})^{\beta-1},
  \label{eq:worker-return}
\end{equation}
falling with the takeover through $1-a$, falling with congestion
through $n^{\beta-1}$, rising with the price. Capital earns
$\beta\,\Pi(\mathbf{r})\,a(\mathbf{r})\,n(\mathbf{r})^{\beta-1}$ per
location. Human and capital payments sum to $\beta V(\mathbf{r})$; the
residual $(1-\beta)V(\mathbf{r})$ is the congestion rent to the fixed
local factor.

\paragraph{Occupation value and equilibrium.}
A worker carries the whole bundle and earns the bundle-integrated
return,
\begin{equation}
  W_o \;=\; \int b_o(\mathbf{r})\, \rho(\mathbf{r})\, d\mathbf{r},
  \label{eq:occ-value}
\end{equation}
the marginal, congestion- and takeover-adjusted counterpart of the gross
bundle price $w_o$ of \eqref{eq:wage-bundle}. Occupation values are
coupled, because a location's density sums over every occupation that
covers it, so occupations that overlap in task space press on one
another's returns. Labor re-sorts from the pre-technology distribution
$L_o^{0}$, drawn toward value and damped by bundle distance $d(o,o')$ in
the sense of \textcite{storck2026mobility}. The allocation is a logit
over value net of the mobility cost from origin,
\begin{equation}
  L_o
  \;=\;
  L^{\text{tot}}\,
  \frac{L_o^{0}\,\exp\!\bigl((W_o - c\, d(o,o^{0}))/\kappa\bigr)}
       {\sum_{o'} L_{o'}^{0}\,\exp\!\bigl((W_{o'} - c\, d(o',o^{0}))/\kappa\bigr)},
  \label{eq:equilibrium}
\end{equation}
a fixed point because $W_o$ depends on $n(\mathbf{r})$, which depends on
$\{L_o\}$. The distance term bites in a static solution because the
re-sorting is measured from an origin, not from nowhere. Total labor is
fixed at $L^{\text{tot}}$; letting it respond to aggregate value is a
later extension.

\paragraph{The labor share.}
Total labor payment is
$\int \beta\,\Pi(\mathbf{r})(1 - a(\mathbf{r}))\,n(\mathbf{r})^{\beta}\,
d\mathbf{r}$ and total factor payment to human and capital is
$\int \beta\,\Pi(\mathbf{r})\,n(\mathbf{r})^{\beta}\, d\mathbf{r}$, so
the labor share of factor income is
\begin{equation}
  \Lambda
  \;=\;
  \frac{\int \Pi(\mathbf{r})\,\bigl(1 - a(\mathbf{r})\bigr)\,
        n(\mathbf{r})^{\beta}\, d\mathbf{r}}
       {\int \Pi(\mathbf{r})\, n(\mathbf{r})^{\beta}\, d\mathbf{r}}.
  \label{eq:labor-share-two-factor}
\end{equation}
The congestion exponent $\beta$ cancels: it governs the spatial sort but
not the share. The share is the price- and density-weighted mean of
$1 - a(\mathbf{r})$. With no automation it is one; the field lowers it
by raising $a$ where \eqref{eq:regime-shift} is crossed. Solving
\eqref{eq:equilibrium} once with the original bundles and once with the
rewritten bundles gives two equilibria, and the difference in $\Lambda$
is the effect of the technology. The displacement of labor onto fewer
tasks then appears not at the automated places, where $a$ is high and
density falls, but at their destinations, where the re-sorting raises
$n$ and congestion bites. The bunching of \textcite{acemoglu2018race} is
recovered as an equilibrium effect of the sort rather than assumed
locally, and whether the share falls hardest in priced or unpriced
directions is decided by where $a$ rises against the price field.

\subsection{Augmentation and the Character of the Technology}
\label{sec:character}

The augmenting part $(1 - s_K)\,\phi_K$ does not take the task. It lowers
the labor the worker needs to produce it. The human labor requirement at
$\mathbf{r}$ is scaled by
\begin{equation}
  \frac{1}{1 + (1 - s_K)\,\phi_K(\mathbf{r})},
  \label{eq:aug-factor}
\end{equation}
so the worker covers the same output with less labor and the output of
the location is unchanged. The labor this releases does not leave the
model. It enters the mobile pool and re-sorts through the allocation
\eqref{eq:equilibrium}, alongside the labor freed by replacement. Where
it lands is set by occupation value and mobility cost, not by the
location it left.

The two parts differ in three ways. \emph{Completeness.} Replacement
takes the whole priced contribution of the crossed micro-tasks, more of
the bundle the higher $s_K$ is and sooner the lower $R$ is, because both
widen the region where \eqref{eq:regime-shift} holds. Augmentation never
takes the task; it scales the requirement by a bounded factor and leaves
the work with the worker. \emph{Feedback.} Replacement is self-damping in
its own direction: as priced mass passes to capital the location's
average price falls, the threshold $R/\Pi$ rises, and further takeover
slows, the brake of \textcite{acemoglu2018race} inherited through the
price field. Augmentation has no such brake, because it moves no priced
mass to capital and does not change $\Pi$. \emph{Destination.}
Replacement seeds reinstatement on the gradient ring, so part of the
displaced mass returns near where it was lost, gated by the deficit.
Augmentation seeds nothing: reinstatement is a fraction of displaced
mass, and the augmenting end displaces nothing, so $M = 0$ and
$\iota_o = 0$ there. The freed labor only re-sorts, and its outcome is
set entirely by the priced locations it can reach. The augmenting end is
a statement about destinations.

The main model takes $s_K$ near one, where capital replaces and
reinstatement follows. The pure augmenting end $s_K = 0$ is the contrast:
no takeover, no seeding, only labor freed and re-sorted.

Two bounds on this treatment. The character $s_K$ is a scalar, uniform
over the reach of the technology. A technology that replaces in one
direction and assists in another is a direction-dependent character,
$s_K(\mathbf{r})$ as a field, which we leave as an extension. The
augmenting end is also where the model most strains its static closure:
separating the labor present at a location from the labor it requires
decouples congestion from output, which the worker layer of
Section~\ref{sec:equilibrium} does not yet carry. We hold the main
results at $s_K$ near one, where labor present and task-work coincide and
the accounting below is exact, and develop the destination problem of
freed labor in work that joins this framework to occupational mobility
\parencite{storck2026mobility}.

%% ============================================================
\section{Calibration}
\label{sec:calibration}
%% ============================================================

\todo{Kalibrering av teknologiparametrar mot empiriska
exponeringsdata (Gmyrek et al. 2025). p\_K fixeras från G3-4-centroid,
z\_K och A\_K kalibreras mot gradient-korrelation och operated task
mass. Jämförelse med Paper 1 Case 2-parametrar. Lägg till kalibrering
av attachment-skalan $\ell$ mot observerade centroidförskjutningar
$\Delta\mu_o$ 2019--2024, och av prisfältets koefficienter $m_0$--$m_5$
direkt från Paper 1. Notera att $\phi_K$ är en generell storhet och att
Gmyrek bara är en realisering: SML-måttet i
\textcite{brynjolfsson2018mlworkforce} ligger redan på task-nivå med
O*NET-importance-vikter och kan tjäna som ett jämförande exponeringsmått
vid sidan av Gmyrek, möjligt just för att talen är tillgängliga.}

%% ============================================================
\section{Illustrative Cases}
\label{sec:cases}
%% ============================================================

\todo{Kondenserad version av Case 1 och Case 2 från Paper 1,
med kalibrerade parametrar. Fokus på den spatiala tröskelmekanismen
och känslighetsanalysen över z\_K. Lägg till en figur över $u(\mathbf{r})$,
kartan över obundna uppgifter, färglagd med $\Pi$.}

%% ============================================================
\section{Empirical Analysis}
\label{sec:empirics}
%% ============================================================

\todo{DiD-specifikation, event study, robusthetscheckar.
Baserat på notebook 9-resultaten. Lägg eventuellt till testet av
buntomritning: projicera $\Delta\mu_o$ på $\nabla\Pi$ och pröva mot
observerade löneförändringar 2019--2024, samma 741 yrken.}

%% ============================================================
\section{Discussion}
\label{sec:discuss}
%% ============================================================

The central mechanism is the directional pricing of work. Reinstatement
is good for employment wherever it lands, because it returns task mass
to labor. It is good for wages only where the price field rewards the
work it returns. New tasks seeded into the analytical directions, where
depth is priced, can raise both employment and the wage. The same new
tasks seeded into the service and manual directions raise employment
but not the wage, because depth is unpriced there. A reinstatement
elasticity that looks protective in aggregate can therefore leave
manual workers with more work and no more pay. This is a distributional
statement that a one-dimensional model cannot make, because it has no
direction in which to place the new work.

Default automation and default reinstatement pull the wage in opposite
directions, and the price field decides the contest. The first term of
\eqref{eq:wage-change} strips the dear tasks, since the operated
threshold \eqref{eq:regime-shift} is crossed first where $\Pi$ is high.
The second term refills with tasks that the deficit gate admits, which
are mostly cheap. Left to run, the process bleeds priced work and
returns unpriced work, the bundle form of the Turing trap. The same
price channel supplies its own brake. As a bundle loses priced mass to
capital, its average price falls, which raises the threshold $R/\Pi$ for
the work that remains and slows further takeover in that direction. This
is the self-correction of \textcite{acemoglu2018race}, where a wave of
automation lowers the effective cost of the labor that is left and
discourages further automation, here inherited through the price field
rather than assumed.

The displacement falls within occupations, not between regions, because
it lives on the interior margin $a(\mathbf{r})$ rather than on a
boundary in space. A bundle spans a range of locations, and the
operated share varies across them, so some of an occupation's own tasks
cross the margin while others do not. This is why the occupation is the
wrong unit and the task is the right one. The point has an empirical
counterpart: the automatability of work varies more within occupations
than between them \parencite{brynjolfsson2018mlworkforce}, which is the
same statement as a bundle spanning a range of $\phi_K$. The within
occupation margin is therefore not a modeling convenience but the level
at which the displacement actually occurs.

Exposure and incidence are distinct, and the model keeps them apart.
The field $\phi_K$ is technical exposure, what capital can do at a
location. The operated share $a(\mathbf{r})$ is economic incidence, what
capital is paid to do there, set by exposure against price through
\eqref{eq:regime-shift}. The two need not coincide. A weak correlation
between exposure and wages does not weaken the trap, because the trap is
a claim about where operation is profitable given exposure, not about
where exposure happens to fall. Directed implementation concentrates on
the work that is both exposed and dear
\parencite{brynjolfsson2018mlworkforce}, which is exactly the high-$\Pi$
high-$\phi_K$ region where \eqref{eq:regime-shift} is crossed first.
Reading exposure alone would miss this; the price field is what turns
exposure into incidence.

The deficit gate is also the boundary between two outcomes that the
debate often conflates. Where attachment capacity $C(\mathbf{r})$ is
positive, the bundle is redrawn and the occupation is redefined. Where
$C(\mathbf{r})$ collapses, the seeded mass stays unbound, and where it
clusters it marks a candidate occupation that no existing bundle can
hold. Redefinition and the birth of a new occupation are then the two
sides of the same gate, separated by whether any occupation already
holds the priced capability the new work requires.

Several limitations bound these claims. The worker layer is closed as a
static equilibrium, so the results are comparative statics between two
rest points rather than a path in time: the adjustment dynamics and the
response of total labor supply to aggregate value are left for later
work. The output value of a place is held fixed under the
human-capital split, which keeps the framework partial and carries no
channel for goods prices or growth. New tasks enter as added weight on
existing locations rather than as genuinely new points on the disk,
which understates the possibility that automation opens directions the
geometry does not yet contain. The cluster weights $v_k$ and the
acquisition scale $\ell$ are identified from the cross-section and held
fixed, and the wage test of the displacement channel, projecting the
centroid shift $\Delta\mu_o$ on $\nabla\Pi$ against observed wage
changes, is the sharpest available check on whether the takeover
framing is right.

\todo{Relatera till Autor \& Thompson (2025) när referensen
verifierats; placera i stycket om riktningsberoende prissättning eller
gränszonen.}

%% ============================================================
\printbibliography
%% ============================================================

\todo{Bibposter att lägga till och verifiera i references.bib:
acemoglu2018race (Acemoglu \& Restrepo, ``The Race between Man and
Machine: Implications of Technology for Growth, Factor Shares, and
Employment'', American Economic Review 108(6):1488--1542, 2018,
DOI 10.1257/aer.20160696 --- verifierad mot fulltext);
brynjolfsson2018mlworkforce (Brynjolfsson, Mitchell \& Rock, ``What Can
Machines Learn and What Does It Mean for Occupations and the Economy?'',
AEA Papers and Proceedings 108:43--47, 2018, DOI 10.1257/pandp.20181019
--- verifierad mot fulltext); brynjolfsson2022turing (Brynjolfsson,
``The Turing Trap: The Promise \& Peril of Human-Like Artificial
Intelligence'', Daedalus 2022 --- ej verifierad mot fulltext, kontrollera
volym, nummer och sidor). Kontrollera även att storck2026geometry,
gmyrek2025 och acemoglu2019automation redan finns. Kontrollera även
nyckeln storck2026mobility, som arbetarlagret nu använder för
buntavståndet $d(o,o')$: bekräfta att den pekar på Paper 2 om
yrkesrörlighet och att nyckelnamnet stämmer.}

\end{document}